\documentclass[10pt]{article}
\usepackage[T1]{fontenc}
\usepackage[utf8]{inputenc}
\usepackage[french]{babel}
\usepackage{amsmath,amssymb,mathrsfs,mathpazo}
\usepackage{geometry}
\geometry{a4paper,margin=2cm}
\usepackage{graphicx,array,tabularx,tabularray,colortbl,multirow,multicol,booktabs,enumitem,pifont,ulem,fancybox,diagbox,cancel}
\usepackage{pstricks,pstricks-add,pst-plot,pst-tree,pst-node,pst-text,pst-func,pst-math,pst-bspline,pst-3dplot,pst-eucl,pst-solides3d}
\let\pscalculate\undefined
\usepackage{pst-bezier}
\usepackage{tikz,pgfplots,tkz-tab}
\usetikzlibrary{arrows,arrows.meta,patterns,calc,positioning,trees,decorations.markings,shapes.geometric,shapes.misc,shapes.arrows,intersections,decorations.pathmorphing,decorations.pathreplacing,backgrounds,shadows,babel}
\pgfplotsset{compat=1.18}
\usepgfplotslibrary{fillbetween}
\usepackage[np]{numprint}
\usepackage{listings,makecell,pas-tableur,twoopt}
\let\barre\undefined
\usepackage[pstricks]{tablvar}
\definecolor{darkgreen}{RGB}{0,100,0}
\providecommand{\red}{\color{red}}
\providecommand{\blue}{\color{blue}}
\providecommand{\green}{\color{green}}
\providecommand{\black}{\color{black}}
\providecommand{\white}{\color{white}}
\providecommand{\gray}{\color{gray}}
\providecommand{\cyan}{\color{cyan}}
\providecommand{\magenta}{\color{magenta}}
\providecommand{\yellow}{\color{yellow}}
\providecommand{\violet}{\color{violet}}
\providecommand{\orange}{\color{orange}}
\providecommand{\euro}{\texteuro}
\providecommand{\EUR}{\texteuro}
\providecommand{\ieme}{\textsuperscript{e}}
\providecommand{\degres}{\ensuremath{^\circ}}
\providecommand{\degree}{\ensuremath{^\circ}}
\providecommand{\vv}[1]{\overrightarrow{#1}}
\providecommand{\Coord}[1]{#1}
\providecommand{\decosix}{}
\providecommand{\starredbullet}{\ensuremath{\bullet}}
\pagestyle{empty}
\setlength{\parindent}{0pt}
\usepackage{adjustbox}\DeclareUnicodeCharacter{2236}{:}\DeclareUnicodeCharacter{2190}{\ensuremath{\leftarrow}}
\providecommand{\R}{\mathbb{R}}
\providecommand{\N}{\mathbb{N}}
\providecommand{\Z}{\mathbb{Z}}
\providecommand{\Q}{\mathbb{Q}}
\providecommand{\D}{\mathbb{D}}
\providecommand{\C}{\mathbb{C}}
\providecommand{\e}{{\mathrm{e}}}
\providecommand{\vect}[1]{\overrightarrow{#1}}
\providecommand{\vectt}[1]{\overrightarrow{\mathrm{#1}}}
\providecommand{\vvt}[1]{\overrightarrow{\mathrm{#1}}}
\providecommand{\np}[2][]{{#2~#1}}
\providecommand{\vv}[1]{{\overrightarrow{#1}}}
\renewcommand{\d}{{\mathrm{d}}}
\providecommand{\degres}{{^\circ}}
\providecommand{\degre}{{^\circ}}
\providecommand{\blue}{\color{blue}}
\providecommand{\red}{\color{red}}
\providecommand{\enskip}{\quad}
\providecommand{\llbracket}{\lbrack\!\lbrack}
\providecommand{\rrbracket}{\rbrack\!\rbrack}
\providecommand{\euro}{{\text{€}}}
\providecommand{\no}{{\text{n°}}}
\providecommand{\interval}[2]{{[#1;#2]}}
\providecommand{\intervFO}[2]{{[#1;#2[}}
\providecommand{\dotfill}{\ldots}

\providecommand{\ds}{\displaystyle}
\providecommand{\textding}[1]{#1}
\providecommand{\Oij}{\(\mathrm{O};\vec{\imath},\vec{\jmath}\)}
\providecommand{\Oijk}{\(\mathrm{O};\vec{\imath},\vec{\jmath},\vec{k}\)}
\providecommand{\Ouv}{\(\mathrm{O};\vec{u},\vec{v}\)}
\providecommand{\intervFF}[2]{\ensuremath{[#1;#2]}}
\renewcommand{\euro}{\ifmmode\text{\texteuro}\else\texteuro\fi}
\DeclareUnicodeCharacter{00B0}{\ifmmode^{\circ}\else\textdegree\fi}
\DeclareUnicodeCharacter{20AC}{\ifmmode\text{\texteuro}\else\texteuro\fi}
% Énoncés : Éducation nationale. Transcription et corrigés : APMEP, Denis Vergès et contributeurs. Découpage et classement : CoopMaths, Sébastien Lozano, Éric Elter et contributeurs. Adaptation de la présentation : Maths Explorer.
% CC BY-SA : https://coopmaths.fr/www/about/
% https://coopmaths.fr/alea/static/bac/2024/tex/bac_2024_05_sujet1_ameriquenord_1.tex
\begin{document}
\noindent\textit{APMEP / CoopMaths. CC BY-SA. Présentation adaptée par Maths Explorer.}\par\medskip
\medskip

Un jeu vidéo récompense par un objet tiré au sort les joueurs ayant remporté un défi. L'objet tiré peut être \og commun \fg{} ou \og rare \fg. Deux types d'objets communs ou rares sont disponibles, des épées et des boucliers.

Les concepteurs du jeu vidéo ont prévu que :

\begin{itemize}
\item la probabilité de tirer un objet rare est de 7\,\% ;
\item si on tire un objet rare, la probabilité que ce soit une épée est de 80\,\% ;
\item si on tire un objet commun, la probabilité que ce soit une épée est de 40\,\%.
\end{itemize}

\emph{Les parties \textbf{\rm A} et \textbf{\rm B} sont indépendantes.}

\bigskip

\textbf{Partie A}

\medskip

Un joueur vient de remporter un défi et tire au sort un objet. On note :

\begin{itemize}[label=$\bullet~$]
\item $R$ l'évènement \og le joueur tire un objet rare\fg{} ;
\item $E$ l'évènement \og le joueur tire une épée \fg{} ;
\item $\overline{R}$ et $\overline{E}$ les évènements contraires des évènements $R$ et $E$.
\end{itemize}

\medskip

\begin{enumerate}
\item Dresser un arbre pondéré modélisant la situation, puis calculer $P(R \cap E)$.
\item Calculer la probabilité de tirer une épée.
\item Le joueur a tiré une épée. Déterminer la probabilité que ce soit un objet rare. Arrondir
le résultat au millième.
\end{enumerate}

\bigskip

\textbf{Partie B}

\medskip

Un joueur remporte $30$ défis.

On note $X$ la variable aléatoire correspondant au nombre d'objets rares que le joueur obtient après avoir remporté $30$ défis. Les tirages successifs sont considérés comme indépendants.

\medskip

\begin{enumerate}
\item Déterminer, en justifiant, la loi de probabilité suivie par la variable aléatoire $X$. Préciser ses paramètres, ainsi que son espérance.
\item Déterminer $P(X < 6)$. Arrondir le résultat au millième.
\item Déterminer la plus grande valeur de $k$ telle que $P(X \geqslant k) \geqslant 0,5$. Interpréter le résultat dans le contexte de l'exercice.
\item Les développeurs du jeu vidéo veulent proposer aux joueurs d'acheter un \og ticket
d'or \fg{} qui permet de tirer $N$ objets. La probabilité de tirer un objet rare reste de 7\,\%.

Les développeurs aimeraient qu'en achetant un ticket d'or, la probabilité qu'un joueur obtienne au moins un objet rare lors de ces $N$ tirages soit supérieure ou égale à $0,95$.

Déterminer le nombre minimum d'objets à tirer pour atteindre cet objectif. On veillera à détailler la démarche mise en œuvre.
\end{enumerate}

\end{document}
